\documentclass[11pt,a4paper]{article}
\usepackage[utf8]{inputenc}
\usepackage[T1]{fontenc}
\usepackage[french]{babel}
\usepackage{lmodern}
\usepackage[margin=2.5cm,headheight=15pt]{geometry}
\usepackage{graphicx,xcolor,fancyhdr,titlesec,enumitem}
\usepackage{amsmath,amssymb,booktabs,caption,float}
\usepackage[most]{tcolorbox}
\usepackage{listings}
\usepackage[hidelinks]{hyperref}

% Palette
\definecolor{principal}{HTML}{1F3A5F}
\definecolor{accent}{HTML}{2E86AB}
\definecolor{fondcode}{HTML}{F5F7FA}
\definecolor{gris}{HTML}{6B7280}

% Titres
\titleformat{\section}{\Large\bfseries\color{principal}}{\thesection.}{0.6em}{}[\vspace{-0.4em}\color{accent}\titlerule]
\titleformat{\subsection}{\large\bfseries\color{accent}}{\thesubsection}{0.6em}{}
\titlespacing*{\section}{0pt}{1.6em}{0.8em}

% En-tête / pied de page
\pagestyle{fancy}
\fancyhf{}
\lhead{\small\color{gris}Développement et déploiement d'applications intelligentes}
\rhead{\small\color{gris}Othmane Moussawi}
\cfoot{\small\color{gris}\thepage}
\renewcommand{\headrulewidth}{0.4pt}

% Encadrés
\newtcolorbox{enonce}{colback=fondcode,colframe=accent,boxrule=0.6pt,
  leftrule=3pt,arc=2pt,title={Énoncé},fonttitle=\bfseries,coltitle=white,colbacktitle=accent}
\newtcolorbox{reponse}{colback=white,colframe=principal,boxrule=0.8pt,
  arc=2pt,title={Réponse},fonttitle=\bfseries,coltitle=white,colbacktitle=principal}

% Code
\lstset{basicstyle=\ttfamily\footnotesize,backgroundcolor=\color{fondcode},
  frame=single,rulecolor=\color{gris},breaklines=true,
  keywordstyle=\color{principal}\bfseries,
  commentstyle=\color{gris}\itshape,showstringspaces=false}

\begin{document}

% ---------- PAGE DE GARDE ----------
\begin{titlepage}
  \centering
  \includegraphics[height=2.4cm]{logo_emsi}\par
  \vspace{2.5cm}
  {\color{gris}\large\textsc{Compte rendu de TP}\par}
  \vspace{0.6cm}
  {\color{principal}\rule{\textwidth}{1.2pt}\par}
  \vspace{0.5cm}
  {\Huge\bfseries\color{principal}Entraîner, évaluer et sauvegarder\\un modèle de sentiment\par}
  \vspace{0.5cm}
  {\color{principal}\rule{\textwidth}{1.2pt}\par}
  \vspace{0.8cm}
  {\Large Développement et déploiement d'applications intelligentes\par}
  \vfill
  {\large\textbf{Réalisé par}\par}
  {\Large Othmane Moussawi\par}
  {\large en binôme avec Ahmed Rayane Ramzi\par}
  \vspace{0.6cm}
  {\large Enseignant : A. Larhlimi\par}
  \vspace{1.5cm}
  {\large Année scolaire : 2026--2027\par}
  {\large 14 octobre 2026\par}
\end{titlepage}
\thispagestyle{empty}
\setcounter{page}{0}
\newpage

\tableofcontents
\newpage

\section{Introduction}

Ce TP transforme le cadrage du TP1 en un modèle concret. Le cas 3 (avis clients négatifs) relève du machine learning : il faut donc entraîner un classifieur de sentiment reproductible, l'évaluer correctement sur des données jamais vues, le sauvegarder pour l'interface du TP3, et comprendre ce qu'il se trompe à corriger. L'objectif n'est pas un modèle optimal, mais un modèle dont on sait ce qu'il vaut et pourquoi il se trompe.

\section{Installation du kit}

\subsection{Démarche}

Le kit \texttt{tp2-semaine2-kit} est installé dans le dépôt par le script \texttt{appliquer\_tp2.py}, d'abord en simulation, puis pour de bon. Le script copie \texttt{src/}, \texttt{tests/test\_model.py} et le notebook, et ajoute \texttt{!models/metrics.json} au \texttt{.gitignore}.

\subsection{Réalisation}

Contrôle C1, \texttt{pytest -q} juste après l'installation : \textbf{3 failed, 37 passed, 2 skipped}. L'énoncé attend 2 failed, 38 passed, 2 skipped. Les deux échecs attendus réclament \texttt{models/metrics.json}, absent avant l'étape C. Le troisième échec est le test d'historique Git (un seul auteur de commits), qui n'est pas lié au modèle.

\subsection{Interprétation}

Le surplus d'échec vient du dépôt, non du kit : le test d'historique Git exige deux auteurs, ce que la séance ne peut pas garantir tant que le second membre n'a pas poussé.

\section{Jeu de démonstration}

\subsection{Démarche}

Le script \texttt{src/make\_demo.py} génère un jeu d'avis en français, déterministe grâce à une graine fixe, composé d'avis nettement positifs, nettement négatifs, mitigés et neutres.

\subsection{Réalisation}

\begin{verbatim}
python src/make_demo.py
\end{verbatim}

Sortie obtenue : 290 avis écrits dans \texttt{data/demo\_reviews.csv}, dont 145 positifs et 145 négatifs, avec 100 avis mitigés et 30 neutres. Le contrôle C2 est conforme.

\subsection{Interprétation}

Les avis mitigés ont une structure fixe : une première moitié d'un sentiment, un retournement (« mais finalement », « sauf que ensuite »), puis une seconde moitié de sentiment opposé. L'étiquette suit la seconde moitié. Cette structure annonce les erreurs de la partie 5.

\section{Entraînement}

\subsection{Démarche}

Le découpage est fait \emph{avant} toute vectorisation, avec \texttt{test\_size=0.2}, une stratification sur l'étiquette et \texttt{random\_state=42}. Vectoriser d'abord laisserait le vocabulaire et les poids du TF-IDF connaître le jeu de test. Le \texttt{Pipeline} enferme le vectoriseur (\texttt{ngram\_range=(1, 2)}, \texttt{min\_df=2}, \texttt{sublinear\_tf=True}) et la régression logistique (\texttt{class\_weight="balanced"}) dans un seul objet, sauvegardé en entier.

\subsection{Réalisation}

\begin{verbatim}
python src/train.py
\end{verbatim}

Sortie obtenue :
\begin{verbatim}
Source : demo, 290 avis (145 négatifs)
232 avis d'entraînement, 58 de test
Vocabulaire : 316 termes
F1 macro sur le jeu de test : 0.879 (seuil 0.85)
\end{verbatim}

Les fichiers produits sont \texttt{models/sentiment.joblib} (non versionné) et \texttt{models/metrics.json} (versionné). Le contrôle C3 est conforme, et seul \texttt{metrics.json} apparaît dans \texttt{git status}.

\subsection{Interprétation}

Ce qui se reconstruit par une commande n'entre pas dans Git. Seule la mesure, \texttt{metrics.json}, est une preuve qui ne se reconstruit pas sans réentraîner : c'est pour cette raison que la CI peut vérifier le seuil sans jamais posséder le modèle.

\section{Évaluation et analyse des erreurs}

\subsection{Démarche}

\texttt{src/evaluate.py} reprend exactement le même découpage (même graine, mêmes proportions) pour produire le rapport de classification et la matrice de confusion.

\subsection{Réalisation}

\begin{verbatim}
python src/evaluate.py
\end{verbatim}

\begin{table}[H]
\centering
\begin{tabular}{lrr}
\toprule
 & \textbf{prédit négatif} & \textbf{prédit positif} \\
\midrule
\textbf{vérité négative} & 27 & 2 \\
\textbf{vérité positive} & 5 & 24 \\
\bottomrule
\end{tabular}
\caption{Matrice de confusion sur les 58 avis de test (jeu de démonstration).}
\end{table}

Le script affiche 7 erreurs sur 58. Le contrôle C4 est conforme : toutes les erreurs portent sur des avis mitigés. Trois exemples relevés :

\begin{itemize}[leftmargin=*]
  \item vérité positive, prédit négatif (confiance 0,53) : « Taille trop petite et retour payant au déballage, sauf que ensuite produit conforme à la description, très satisfait. »
  \item vérité négative, prédit positif (confiance 0,56) : « Service impeccable et livraison en avance sur le papier, sauf que ensuite produit cassé à l'arrivée. »
  \item vérité négative, prédit positif (confiance 0,69) : « Service impeccable et livraison en avance les premiers jours, avant que notice incompréhensible et pièces manquantes. »
\end{itemize}

\subsection{Interprétation}

Les deux faux positifs sont le cas grave : un client mécontent traité en retard. Le mécanisme est celui d'un sac de mots, qui ne tient pas compte de l'ordre : « impeccable … puis cassé » et « cassé … puis impeccable » donnent un vecteur presque identique. Le modèle hésite d'ailleurs sur ces cas, avec des confiances entre 0,53 et 0,69. Corriger ce défaut demande une représentation sensible à l'ordre, comme des embeddings (semaine 9).

\section{Le vrai jeu : Allociné}

\subsection{Démarche}

Le même code est lancé sur 20 000 critiques de films étiquetées, téléchargées depuis Hugging Face, avec le même découpage. Le jeu est près de soixante-dix fois plus grand que celui de démonstration.

\subsection{Réalisation}

\begin{verbatim}
python src/train.py --allocine
python -m pytest -q
\end{verbatim}

Sorties obtenues :
\begin{verbatim}
Source : allocine, 20000 avis (10204 négatifs)
16000 avis d'entraînement, 4000 de test
Vocabulaire : 149529 termes
F1 macro sur le jeu de test : 0.920 (seuil 0.85)
\end{verbatim}

\begin{table}[H]
\centering
\begin{tabular}{lrr}
\toprule
 & \textbf{Démonstration} & \textbf{Allociné} \\
\midrule
Avis & 290 & 20 000 \\
Vocabulaire & 316 termes & 149 529 termes \\
F1 macro & 0,879 & 0,920 \\
\bottomrule
\end{tabular}
\caption{Comparaison des deux jeux.}
\end{table}

Matrice de confusion sur 4 000 avis de test Allociné : 1 860 négatifs bien classés, 181 négatifs classés positifs, 137 positifs classés négatifs, 1 822 positifs bien classés (318 erreurs, soit 7,95 %).

Avec \texttt{pytest -q}, on obtient \textbf{41 passed, 1 failed} : le test en échec est l'historique Git, qui attend le commit du second membre. Le contrôle C5 attend 42 passed après ce commit.

\subsection{Interprétation}

Le modèle n'a pas changé : ce sont les données qui ont changé. La mesure devient crédible, calculée sur 4 000 avis au lieu de 58. Sur Allociné, les négatifs manqués sont plus nombreux que les positifs manqués, à l'inverse de la démonstration.

\begin{reponse}
Modèle sauvegardé avec F1 macro 0,920 sur 4 000 avis. Erreurs principales : le sac de mots est insensible à l'ordre des phrases, ce qui pénalise les avis mitigés comme les textes qui se contredisent.
\end{reponse}

\section*{Conclusion}

Le modèle est entraîné, évalué sur un jeu de test jamais vu, et sa mesure est versionnée. Le TP a aussi montré que les performances dépendent autant des données que du code : le même pipeline passe de 0,879 à 0,920 selon le jeu utilisé.

\end{document}
